\section{Introducción}
\label{sec:intro}

La inflación cósmica predice que, además de las fluctuaciones de densidad que sembraron las galaxias, el universo primitivo generó un fondo de ondas gravitacionales. Su amplitud está fijada por la escala de energía a la que ocurrió la inflación: en relatividad general (RG), el espectro de potencia de esas ondas vale, a orden más bajo, $2H^2/(\pi^2\MP^2)$ evaluado en el momento en que cada modo sale del horizonte, donde $H$ es el parámetro de Hubble y $\MP$ la masa de Planck reducida. Por eso el espectro tensorial primordial es una de las pocas ventanas observacionales hacia la física de la inflación, y por eso las cotas actuales sobre la relación tensor-escalar $r$ (Planck, BICEP/Keck) restringen tanto los modelos.

La gravedad unimodular (UG, por \emph{unimodular gravity}) es una teoría de la gravitación que, en su versión más simple, reproduce clásicamente las ecuaciones de la RG, pero trata la constante cosmológica de un modo distinto: en lugar de ser un parámetro de la acción, aparece como una constante de integración, y la energía del vacío deja de gravitar directamente. Esa es su motivación histórica, ligada al problema de la constante cosmológica. Lo que interesa a este trabajo es una variante que la propia estructura de la teoría permite: si se admite que la energía-impulso de la materia no se conserva, aparece un escalar $Q(x)$, llamado a veces \emph{término de difusión}, que actúa sobre el fondo cosmológico como un fluido con ecuación de estado $w=-1$ pero con densidad variable. Esta variante fue explorada como modelo de inflación por G.~León y colaboradores \cite{leon22,piccirilli23}, y su motivación física (modificaciones no unitarias de la mecánica cuántica, o discretización del espacio-tiempo) se discute en \cite{josset17,perez19}.

\subsection*{La pregunta}

Con $Q$ presente, hay dos preguntas naturales sobre el espectro tensorial primordial, y son las que ordenan todo el informe:
\begin{enumerate}[label=(\alph*),nosep]
\item ¿Cambia el fondo cosmológico durante la inflación? En particular, ¿cómo entra $Q$ en las ecuaciones de Friedmann y en los parámetros de flujo de Hubble?
\item ¿Cambia la ecuación de los modos tensoriales y, por tanto, el espectro $P_T(k)$? Es decir, ¿el vínculo unimodular, o la no conservación de la materia, alteran la propagación de las ondas gravitacionales, o solo alteran el fondo sobre el que se propagan?
\end{enumerate}
La literatura que se pudo leer (sección~\ref{sec:ug-literatura}) da una respuesta cualitativa a la segunda pregunta: ningún trabajo encuentra diferencias en el sector tensorial atribuibles al vínculo, y las discrepancias publicadas están en el sector escalar. Pero también da una señal de que la pregunta está menos cerrada de lo que parece. Los trabajos de León \cite{leon22} derivan el espectro escalar con difusión pero no contienen sector tensorial; y el trabajo de Piccirilli y León \cite{piccirilli23}, que sí da la relación $r=16\,\epsilon_1$, la adopta sin derivarla, remitiendo a \cite{leon22}. Que el sector tensorial coincida con el de RG es, en ese contexto, un resultado que se da por supuesto más de lo que se demuestra, y en la literatura leída no encontré una derivación explícita desde la acción a segundo orden con el vínculo y con difusión.

\subsection*{Qué se hizo}

El trabajo tiene tres partes. La primera es una derivación analítica, hecha a mano, que parte de la formulación de UG con multiplicador de Lagrange y llega, sobre un fondo de Friedmann–Lemaître–Robertson–Walker (FLRW) plano, a las ecuaciones de fondo con $Q$, a la ecuación de los modos tensoriales y a la acción cuadrática, incluyendo el término del vínculo a segundo orden, que no es trivial (secciones~\ref{sec:fondo} y \ref{sec:tensores}). La segunda es un cálculo numérico del espectro $P_T(k)$ en RG y en UG para un inflatón con potencial cuadrático y con el de Starobinsky, integrando modo a modo la ecuación de los modos con condiciones de Bunch–Davies (sección~\ref{sec:numerico}). La tercera es una batería de controles: comparación con el límite estándar de RG, convergencia numérica, y cuantificación de la diferencia entre UG y RG con la demostración de que no es un artefacto numérico (secciones~\ref{sec:resultados} y~\ref{sec:potenciales}). Se agregan la comparación a igual número de e-folds antes del final de la inflación (sección~\ref{sec:potenciales}) un análisis del acoplamiento de $Q$ a la materia, con verificación simbólica (sección~\ref{sec:acoplamiento}), y el cálculo con el modelo de radiación con difusión Q \cite{leon22,piccirilli23}, radiación pura más $Q$ (sección~\ref{sec:leon}).

Conviene ser explícito desde ahora sobre lo que \emph{no} se hizo. No se trata el sector escalar de UG con difusión, de modo que no se calcula $r$ ni se resuelve la discrepancia entre las expresiones de $P_\mathcal{R}$ de \cite{leon22} y \cite{piccirilli23}, que se describe en la sección~\ref{sec:discusion}. De la derivación, solo la ecuación de los tensores a primer orden y algunas identidades del fondo y del acoplamiento de $Q$ se verificaron con álgebra simbólica (sección~\ref{sec:acoplamiento}); la acción a segundo orden sigue verificada solo a mano. Y los resultados numéricos son para un único conjunto de parámetros de $Q$, cuyos valores fijan buena parte de la magnitud de la diferencia entre UG y RG.

\subsection*{Organización}

Las secciones~\ref{sec:rg} y \ref{sec:inflacion} son material elemental y estándar: relatividad general y cosmología homogénea, y la inflación con un campo escalar, incluyendo el cálculo del espectro tensorial en RG, que sirve de referencia para todo lo que sigue. Quien ya domine ese material puede pasar directamente a la sección~\ref{sec:ug}, que presenta la gravedad unimodular y su término de difusión, y a las secciones~\ref{sec:fondo} y~\ref{sec:tensores}, donde está el núcleo analítico. La sección~\ref{sec:espectro} junta la derivación con la cuantización, la sección~\ref{sec:numerico} describe el cálculo numérico, la sección~\ref{sec:resultados} presenta los resultados y los controles, y la sección~\ref{sec:discusion} discute qué se puede concluir y de qué depende. Los apéndices contienen el álgebra más larga, la tabla de procedencia y el detalle de fuentes y de código.

\subsection*{Convenciones}

Se usa la signatura $(-,+,+,+)$, unidades $c=\hbar=1$ y $\kap\equiv 8\pi G=1/\MP^2$, con $\MP$ la masa de Planck reducida. Las fórmulas del texto son dimensionales, con $\MP$ explícita, salvo donde se declara \emph{expresamente} $\MP=1$ (secciones~\ref{sec:numerico} en adelante y el código, que además usa $m=1$); en esos pasajes $\Lambda_0$ se mide en unidades de $\MP^2$, $Q$, $V$ y $\rho$ en unidades de $\MP^4$, y $V_{\rm eff}=V+Q+\Lambda_0\MP^2$ se escribe $V+Q+\Lambda_0$. Los índices latinos $i,j,k$ son espaciales y se suman cuando se repiten; $h_{ij}^2\equiv h_{ij}h_{ij}$. El punto denota derivada respecto del tiempo cósmico $t$; la prima, respecto del tiempo conforme $\eta$, con $\dd t=a\,\dd\eta$, y $\mathcal H\equiv a'/a=aH$. El número de e-folds es $N=\ln(a/a_i)$. El espacio se toma plano.
